% ═══════════════════════════════════════════════════════════════
% MT-08d-ML: Musterlösung Minitest Repaso - Jahreswiederholung
% Spanisch Klasse 7
% ═══════════════════════════════════════════════════════════════
% Dauer: 20 Minuten
% Punkte: 30 BE
% Inhalt: Verkehrsmittel, Wetter, Pretérito, ir a + Infinitiv
% ═══════════════════════════════════════════════════════════════

\documentclass[11pt, a4paper]{article}

\newif\ifswmodus
\swmodusfalse

\usepackage[utf8]{inputenc}
\usepackage[T1]{fontenc}
\usepackage[ngerman]{babel}
\usepackage{helvet}
\renewcommand{\familydefault}{\sfdefault}

\usepackage[margin=2cm]{geometry}
\usepackage{enumitem}
\usepackage{tabularx}
\usepackage{booktabs}

\usepackage{xcolor}
\usepackage{tcolorbox}
\tcbuselibrary{skins}

\usepackage{fancyhdr}
\usepackage{lastpage}
\usepackage{needspace}

% FARBEN
\definecolor{spanisch-color}{HTML}{c41e3a}
\definecolor{grau-color}{HTML}{888888}
\definecolor{hellgrau-color}{HTML}{f5f5f5}
\definecolor{loesung-color}{HTML}{c62828}

\definecolor{sw-dark}{HTML}{000000}
\definecolor{sw-gray}{HTML}{666666}
\definecolor{sw-white}{HTML}{ffffff}

\ifswmodus
    \colorlet{spanisch}{sw-dark}
    \colorlet{grau}{sw-gray}
    \colorlet{hellgrau}{sw-white}
    \colorlet{loesung}{sw-dark}
\else
    \colorlet{spanisch}{spanisch-color}
    \colorlet{grau}{grau-color}
    \colorlet{hellgrau}{hellgrau-color}
    \colorlet{loesung}{loesung-color}
\fi

\newcommand{\fachfarbe}{spanisch}

\newcommand{\thema}{Minitest: Repaso -- MUSTERLÖSUNG}
\newcommand{\fach}{Spanisch}
\newcommand{\klasse}{7}
\newcommand{\maxpunkte}{30}
\newcommand{\zeit}{20 Minuten}
\newcommand{\datum}{\today}

\pagestyle{fancy}
\fancyhf{}
\renewcommand{\headrulewidth}{0.4pt}
\renewcommand{\footrulewidth}{0.4pt}

\lhead{\fach{} -- Klasse \klasse}
\chead{\textbf{MT-08d -- ML}}
\rhead{\datum}

\lfoot{\zeit}
\cfoot{}
\rfoot{Seite \thepage\ von \pageref{LastPage}}

\newcommand{\punktebox}[1]{\hfill\fbox{\small \_/#1}}
\newcommand{\luecke}[1]{\rule{#1}{0.4pt}}
\newcommand{\lsg}[1]{\textcolor{loesung}{\textbf{#1}}}

\newtcolorbox{infobox}{
    colback=hellgrau,
    colframe=\fachfarbe,
    boxrule=1pt,
    arc=0pt,
    left=8pt, right=8pt, top=8pt, bottom=8pt
}

\newenvironment{aufgabe}[1]{%
    \needspace{5\baselineskip}%
    \nopagebreak[4]%
    \vspace{10pt}
    \noindent\textbf{\textcolor{\fachfarbe}{Aufgabe #1}}
    \nopagebreak[4]%
    \vspace{5pt}
    \nopagebreak[4]%
    \begin{list}{}{\leftmargin=0pt}
    \item[]
}{%
    \end{list}
}

\newcommand{\es}[1]{\textcolor{\fachfarbe}{\textbf{#1}}}

\begin{document}

% ═══════════════════════════════════════════════════════════════
% KOPFBEREICH
% ═══════════════════════════════════════════════════════════════

\begin{center}
    {\Large\bfseries\textcolor{\fachfarbe}{\thema}}
\end{center}

\vspace{5pt}

\noindent
\begin{tabularx}{\textwidth}{|X|X|X|}
\hline
\textbf{Name:} \lsg{MUSTERLÖSUNG} &
\textbf{Klasse:} \klasse &
\textbf{Datum:} \luecke{2cm} \\
\hline
\end{tabularx}

\vspace{10pt}

\begin{infobox}
\textbf{Zeit:} \zeit{} | \textbf{Punkte:} \maxpunkte{} BE | \textbf{Hilfsmittel:} keine \\
\textbf{Themen:} Verkehrsmittel, Wetter, Pretérito Indefinido, ir a + Infinitiv
\end{infobox}

\vspace{15pt}

% ═══════════════════════════════════════════════════════════════
% AUFGABE 1: Verkehrsmittel (4 BE)
% ═══════════════════════════════════════════════════════════════

\begin{aufgabe}{1 -- Verkehrsmittel und Orte}
Verbinde das spanische Wort mit der deutschen Bedeutung. \punktebox{4}
\end{aufgabe}

\begin{center}
\begin{tabular}{rcl}
\es{el avión} & \lsg{2} & der Zug \\[0.8em]
\es{el tren} & \lsg{1} & das Flugzeug \\[0.8em]
\es{la playa} & \lsg{4} & das Hotel \\[0.8em]
\es{el hotel} & \lsg{3} & der Strand \\
\end{tabular}
\end{center}

\textit{Lösung: el avión = das Flugzeug, el tren = der Zug, la playa = der Strand, el hotel = das Hotel}

\vspace{15pt}

% ═══════════════════════════════════════════════════════════════
% AUFGABE 2: Wetter (6 BE)
% ═══════════════════════════════════════════════════════════════

\begin{aufgabe}{2 -- Wetter übersetzen}
Übersetze ins Spanische. \punktebox{6}
\end{aufgabe}

\noindent
a) Es ist warm. $\rightarrow$ \lsg{Hace calor.}

\vspace{0.8em}

\noindent
b) Es regnet. $\rightarrow$ \lsg{Llueve.}

\vspace{0.8em}

\noindent
c) Es ist kalt. $\rightarrow$ \lsg{Hace frío.}

\vspace{15pt}

% ═══════════════════════════════════════════════════════════════
% AUFGABE 3: Pretérito Konjugation (8 BE)
% ═══════════════════════════════════════════════════════════════

\begin{aufgabe}{3 -- Pretérito Indefinido konjugieren}
Konjugiere die Verben im Pretérito Indefinido. \punktebox{8}
\end{aufgabe}

\noindent
a) yo / viajar: \lsg{viajé}

\vspace{0.8em}

\noindent
b) tú / visitar: \lsg{visitaste}

\vspace{0.8em}

\noindent
c) él / nadar: \lsg{nadó}

\vspace{0.8em}

\noindent
d) nosotros / bailar: \lsg{bailamos}

\vspace{15pt}

% ═══════════════════════════════════════════════════════════════
% AUFGABE 4: ir a + Infinitiv (6 BE)
% ═══════════════════════════════════════════════════════════════

\begin{aufgabe}{4 -- ir a + Infinitiv}
Ergänze die richtige Form von \textit{ir a + Infinitiv}. \punktebox{6}
\end{aufgabe}

\noindent
a) Yo \lsg{voy a ir} a la playa. \hfill (gehen)

\vspace{0.8em}

\noindent
b) Nosotros \lsg{vamos a viajar} en avión. \hfill (reisen)

\vspace{0.8em}

\noindent
c) Ellos \lsg{van a nadar} en el mar. \hfill (schwimmen)

\vspace{15pt}

% ═══════════════════════════════════════════════════════════════
% AUFGABE 5: Jahreszeiten (6 BE)
% ═══════════════════════════════════════════════════════════════

\begin{aufgabe}{5 -- Jahreszeiten und Wetter}
Ergänze passend. \punktebox{6}
\end{aufgabe}

\noindent
a) En verano \lsg{hace calor y hace sol}.

\vspace{0.8em}

\noindent
b) En invierno \lsg{hace frío y nieva}.

\vspace{0.8em}

\noindent
c) \es{la primavera} = \lsg{der Frühling}

% ═══════════════════════════════════════════════════════════════
% PUNKTEÜBERSICHT
% ═══════════════════════════════════════════════════════════════

\vspace{25pt}
\hrule
\vspace{10pt}

\noindent
\begin{tabularx}{\textwidth}{|X|c|c|c|c|c||c|}
\hline
\textbf{Aufgabe} & 1 & 2 & 3 & 4 & 5 & \textbf{Gesamt} \\
\hline
\textbf{max. Punkte} & 4 & 6 & 8 & 6 & 6 & \textbf{30} \\
\hline
\textbf{erreicht} & & & & & & \\[0.8em]
\hline
\end{tabularx}

\vspace{10pt}

\begin{flushright}
\textbf{Note:} \luecke{2cm}
\end{flushright}

\end{document}
